\section{Discussion}
\label{sec:discussion}

\paragraph{The mixture is a poor labeler and a useful tie-breaker.}
Used on its own, the mixture is the worst row in Table~\ref{tab:p1first}. Its
components mostly separate Recruitment from the other reasons. Adding more
components splits Recruitment into smaller pieces. It does not isolate Safety or
Funding. The embeddings do contain those reasons: a supervised model reaches
0.941 F1 on Safety from the same inputs. The problem is the clustering. It finds
the largest pattern in the data, and that pattern is Recruitment.

The mixture still earns its place once it is demoted. In the first family, using
it as a 20\% voice in self-training raised cross-validated macro F1 from 0.805 to
0.814. In the submitted model we cut that voice to 5\%. Its job is to steady
which unlabeled trials get promoted, not to name their reasons. The quarter
weight on a generated row does the same kind of work. In the final fit the 757
accepted rows carry about half the weight of the 235 real labels, so a confident
guess can help and cannot take over.

\paragraph{Why adding every feature makes the model worse.}
This was the clearest negative result. The embeddings are a compressed copy of
the stopping-reason text, which is the text the labels came from. They already
hold most of what the labels have to say. Criteria, diseases, and drugs describe
what the trial was, not why it ended. Adding one of those blocks means 20 to
2,000 new inputs on 157 real labels. The noise shows up immediately. Eligibility
text drops validation sample-wise F1 from 0.893 to 0.876 in
Table~\ref{tab:p1features}. Fingerprints barely move the cross-validated score,
and 44\% of trials have no matched drug. Start year looks predictive only because
the labeled trials were chosen to be early. The test trials are the late ones,
so a year pattern learned here points the wrong way. Leaving a feature out was a
result, not a shortcut.

\paragraph{A higher validation score was not enough.}
The linear-discriminant ensemble beats the submitted blend by 0.015 on the 78
validation trials and is tied with it under cross-validation. We did not follow
the validation number. Four Funding trials and ten Other trials can move the
macro average by several points, which is why the validation and
cross-validated columns in Table~\ref{tab:p1first} already disagree about which
early model is best. The time-ordered windows, which look more like the later
test years, favor the submitted blend. So does the reported Kaggle score. The
validation score also belongs to a fit that has not yet seen those 78 labels.
The submitted file is retrained with them included.

\paragraph{What we would not do again.}
We would not fit the scaler or the mixture on the test rows. That choice is part
of the 0.89423 submission, and the later inductive fit scored higher. We would
not let generated labels rewrite the class weights, because the model then starts
believing its own class balance. We would not add year, text, or drug features
in the hope that more inputs are safer. The gain that held up was smaller than
that search: freeze the weights, keep both of a neighbor's reasons, and let the
mixture advise instead of decide.
